\thispagestyle{empty}
\begin{center}
{\heiti\fontsize{18}{23}\selectfont 多无人机协同计算的双时间尺度调度\par}
\vspace{8mm}
{\fontsize{12}{17}\selectfont 作者姓名待填写\par}
\vspace{3mm}
{\fontsize{10.5}{15}\selectfont （作者单位全称，省市 邮政编码待填写）\par}
\end{center}
\vspace{5mm}
\noindent\textbf{摘\quad 要：}
针对动态热点下多无人机计算负载失衡及部署中断问题，建立了包含上传、转发、计算和分片同步完成的任务级模型。采用近端策略优化确定大时间尺度部署，采用软演员评论家确定小时间尺度路由，并通过单纯形映射满足分配约束。以未完成父任务的加权驻留量和能耗构造共同代价，按物理时间统一两层折扣，补全跨部署边界的经验。给出了动作可行性、价值表征及复杂度分析，为连续任务服务下的分层调度提供建模与求解方案。

\vspace{3mm}
\noindent\textbf{关键词：}无人机；移动边缘计算；双时间尺度；任务路由；深度强化学习

\vspace{3mm}
\noindent\textbf{中图分类号：}待作者确认\qquad
\textbf{文献标志码：}A\qquad
\textbf{DOI：}待编辑部编定

\vspace{10mm}
\begin{center}
{\bfseries\fontsize{15}{20}\selectfont Two-Timescale Scheduling for Cooperative Multi-UAV Computing\par}
\vspace{4mm}
{\fontsize{11}{16}\selectfont AUTHOR Names to Be Supplied\par}
\vspace{2mm}
{\fontsize{10}{14}\selectfont Full affiliation, City Postal Code, China\par}
\end{center}
\vspace{4mm}
\noindent\textbf{Abstract:}
Dynamic hotspots and deployment interruptions were considered in cooperative unmanned aerial vehicle (UAV) computing. A task-level model was formulated to describe uploading, forwarding, computation, and synchronized branch completion. Proximal policy optimization determined deployment at the large time scale, while soft actor-critic determined routing at the small time scale. A simplex mapping enforced allocation constraints. Weighted unfinished-task residence and energy formed a common cost. Discount factors were aligned with physical time, and transitions spanning deployment interruptions were completed using their actual successor observations. Action feasibility, value representation, and computational complexity were analyzed. The formulation provides a modeling and solution framework for hierarchical scheduling with continuous task service.

\vspace{4mm}
\noindent\textbf{Key words:} unmanned aerial vehicle; mobile edge computing; two time scales; task routing; deep reinforcement learning

\vfill
\noindent{\small 稿件状态：系统模型、决策建模与算法讨论稿；联合训练性能及正式对比实验尚待补充。此说明在具备完整实验证据后删除。}

\vspace{4mm}
\noindent{\small 通信作者：待填写；邮箱：待填写。\par
基金项目：按真实资助情况填写；无资助则删除本项。}
